% Magic squares of squares and elliptic divisibility sequences: the rank-one case.
% Draft, 2026-11-05. Compile with pdflatex (no special packages). Before submission: the author should
% review the whole argument, and double-check bibliographic details such as page ranges.
\documentclass[11pt]{amsart}
\usepackage{amsmath,amssymb,amsthm}
\usepackage[hidelinks]{hyperref}

\newtheorem{theorem}{Theorem}[section]
\newtheorem{proposition}[theorem]{Proposition}
\newtheorem{lemma}[theorem]{Lemma}
\newtheorem{corollary}[theorem]{Corollary}
\theoremstyle{definition}
\newtheorem{definition}[theorem]{Definition}
\theoremstyle{remark}
\newtheorem{remark}[theorem]{Remark}
\newtheorem{example}[theorem]{Example}

\newcommand{\Q}{\mathbb{Q}}
\newcommand{\Z}{\mathbb{Z}}
\DeclareMathOperator{\sqf}{sqf}
\DeclareMathOperator{\rk}{rank}
\hyphenation{Tza-na-kis Ver-zo-bio}

\title[Magic squares of squares: the rank-one case]{Magic squares of squares and elliptic divisibility
sequences: the rank-one case}
\author{Ryan Heard}
\date{\today}
\subjclass[2020]{11D09, 11G05, 11B39}
\keywords{magic square of squares, congruent number curve, elliptic divisibility sequence, primitive divisor}

\begin{document}

\begin{abstract}
Robertson observed that a $3\times3$ magic square of squares gives three points on a congruent number
curve $E_n\colon y^2=x^3-n^2x$ whose $x$-coordinates are in arithmetic progression, and three more such
points on a second congruent number curve. Bremner, Silverman and Tzanakis proved that for squarefree $n$
a subgroup of $E_n(\Q)$ of rank one contains no non-trivial \emph{integral} arithmetic progression. Their
theorem never applies to magic squares of squares: in a primitive one the common differences are
divisible by~$8$, and as a result the points never have integral coordinates. We prove a
rational analogue on the curves $y^2=x^3+Dx$: the $x$-coordinates of three multiples of a non-torsion
point with square $x$-coordinate, with pairwise distinct absolute indices, are never in arithmetic
progression. The proof uses the primitive divisor theorem of Voutier and Yabuta, with Verzobio's
correction. Consequently, in any magic square of squares each of the two families of three points
generates a subgroup of rank at least two. Equivalently, the common difference of either diagonal of the
square has a squarefree part $n$ with $\rk E_n(\Q)\ge2$. By a theorem of Smith, such $n$ form a set of
density zero.
\end{abstract}

\maketitle

\section{Introduction}

A \emph{magic square of squares} is a $3\times3$ array of nine distinct positive integer squares whose
rows, columns and two diagonals all have the same sum. Whether one exists is a well-known open problem,
popularized by LaBar~\cite{LaBar} and by Gardner (see Boyer~\cite{Boyer} for the history, the prizes, and
near-misses such as Bremner's magic square with seven square entries). Robertson~\cite{Robertson} related
the problem to congruent number curves, and Bremner~\cite{Bremner99,Bremner01} studied it through its
sub-configurations and their geometry. Bruin, Thomas and V\'arilly-Alvarado~\cite{BTVA} proved that the
associated surface is algebraically quasi-hyperbolic. Auel and Singer~\cite{AS} have recently described
its algebraic geometry in detail, and Kominers~\cite{Kominers} has studied magic squares of squares with
center~$0$ over number fields. Extensive computer searches have found no example.

Every $3\times3$ magic square of integers can be written in the Lucas form
\begin{equation}\label{eq:lucas}
\begin{pmatrix}
E+a & E-a-b & E+b\\
E-a+b & E & E+a-b\\
E-b & E+a+b & E-a
\end{pmatrix},
\qquad E,a,b\in\Z,
\end{equation}
whose entries are the nine numbers $L_{\alpha,\beta}=E+\alpha a+\beta b$ with
$\alpha,\beta\in\{-1,0,1\}$. The entries are distinct exactly when $ab\neq0$, $a\neq\pm b$,
$a\neq\pm2b$ and $b\neq\pm2a$. The main diagonal is $(E-a,E,E+a)$ and the anti-diagonal is
$(E-b,E,E+b)$.

For fixed $\alpha$, the three entries $L_{\alpha,-1},L_{\alpha,0},L_{\alpha,1}$ form an arithmetic
progression with common difference $b$. In~\eqref{eq:lucas} they are the anti-diagonal ($\alpha=0$) and
the two broken diagonals parallel to it. We call these the three \emph{$b$-progressions}; they partition
the nine entries. Likewise, for fixed $\beta$, the entries $L_{-1,\beta},L_{0,\beta},L_{1,\beta}$ form
the three \emph{$a$-progressions}, with common difference $a$: the main diagonal and its two broken
parallels.

For a nonzero integer $m$, write $\sqf(m)$ for the squarefree part of $|m|$: the squarefree integer
$n\ge1$ with $|m|=nk^2$ for some integer $k$. For a squarefree $n\ge1$, let $E_n$ be the congruent number
curve
\[
E_n\colon\ y^2=x^3-n^2x.
\]
If all nine entries of~\eqref{eq:lucas} are squares, write $|b|=nk^2$ with $n=\sqf(b)$. For each $\alpha$,
the three entries of the $\alpha$-th $b$-progression divided by $k^2$ are the rational squares $x_\alpha-n$,
$x_\alpha$ and $x_\alpha+n$, where $x_\alpha=(E+\alpha a)/k^2$. So $x_\alpha(x_\alpha-n)(x_\alpha+n)$ is a
rational square, and there is a point
\begin{equation}\label{eq:points}
P_\alpha=(x_\alpha,y_\alpha)\in E_n(\Q),\qquad x_{-1}+x_1=2x_0 .
\end{equation}
This is the observation of Robertson~\cite{Robertson}. The $a$-progressions give three points on
$E_{\sqf(a)}$ in the same way.

Motivated by this connection, Bremner, Silverman and Tzanakis proved in~\cite{BST} that if $n$ is
squarefree then a subgroup of $E_n(\Q)$ of rank one contains no non-trivial \emph{integral} arithmetic progression.
Their proof uses delicate height estimates. Their theorem does not reach magic squares of squares. In a
primitive magic square of squares all nine entries are odd and $8$ divides $a$ and $b$
(Remark~\ref{rem:BST}), so $b$ is never squarefree, and the points~\eqref{eq:points} never have integral
coordinates. Our main result is a rational version, for the points that magic squares actually produce.

\begin{theorem}\label{thm:main}
Let $E,a,b$ be integers with $ab\neq0$ such that all nine entries of~\eqref{eq:lucas} are squares of
integers. Then the points $P_{-1},P_0,P_1$ of~\eqref{eq:points} generate a subgroup of $E_{\sqf(b)}(\Q)$
of rank at least two. Likewise, the three points coming from the $a$-progressions generate a subgroup of
$E_{\sqf(a)}(\Q)$ of rank at least two.
\end{theorem}

\begin{corollary}\label{cor:rank}
Let $d$ be the common difference of either diagonal of a magic square of squares. Then $n=\sqf(d)$ is a
congruent number with $\rk E_n(\Q)\ge2$. The set of such $n$ has natural density zero among squarefree
integers. Its elements up to $150$ are $34,41,65,137,138,145$.
\end{corollary}

\begin{example}\label{ex:6}
The curve $E_6$ has rank one; its point $2G$, for the generator $G=(-3,9)$, has $x(2G)=25/4$, which
corresponds to the progression $1,25,49$ of difference $24=6\cdot2^2$. By Corollary~\ref{cor:rank}, no
magic square of squares has a diagonal whose common difference is of the form $6k^2$. In particular,
$(1,25,49)$ is not a diagonal of any magic square of squares.
\end{example}

The key input is Lemma~\ref{lem:noAP}. Let $Q$ be a non-torsion point on an elliptic curve with a
Weierstrass equation with integer coefficients, and suppose that every term of the elliptic divisibility
sequence of $Q$ from the third on has a primitive divisor. Then the $x$-coordinates of three multiples
of $Q$ with pairwise distinct absolute indices are never in arithmetic progression. Its proof is short.
On the curves $y^2=x^3+Dx$ with $D$ fourth-power-free in minimal form, the primitive divisor theorem of
Voutier and Yabuta~\cite{VY}, which appeared after~\cite{BST}, supplies the hypothesis whenever $x(Q)$
is a square (Theorem~\ref{thm:VY}). More
general results on arithmetic progressions among $x$-coordinates of rational points, such as Choi's
additive rigidity theorem~\cite{Choi}, bound the number of points in such progressions in terms of the
Mordell--Weil rank, but do not give this exclusion. As far as we know, Lemma~\ref{lem:noAP} and
Theorem~\ref{thm:main} are new. Both, and the rank statement of Corollary~\ref{cor:rank}, have been
checked in the Lean proof assistant, with Theorem~\ref{thm:VY} and the torsion of $E_n(\Q)$ as explicit
hypotheses (Remark~\ref{rem:lean}).

\section{Congruent number curves and their divisibility sequences}

\begin{lemma}\label{lem:facts}
Let $n\ge1$ be squarefree.
\begin{enumerate}
\item[(a)] $E_n(\Q)_{\mathrm{tors}}=E_n[2]=\{O,(0,0),(n,0),(-n,0)\}$.
\item[(b)] A point $P=(x,y)\in E_n(\Q)$ lies in $2E_n(\Q)$ if and only if $x$, $x-n$ and $x+n$ are all
squares in $\Q$.
\item[(c)] $2E_n(\Q)$ is torsion-free.
\item[(d)] The equation $y^2=x^3-n^2x$ is a global minimal model of $E_n$.
\end{enumerate}
\end{lemma}

\begin{proof}
Part (a) is classical; see~\cite[Ch.~I, \S9]{Ko}. Part (b) is the standard criterion for halving a point
on $y^2=(x-e_1)(x-e_2)(x-e_3)$ with $e_i\in\Q$, applied with $\{e_i\}=\{0,n,-n\}$; see~\cite[Ch.~IV]{Kn}.
For (c), if $T=2R\in2E_n(\Q)$ is torsion then so is $R$, so $R\in E_n[2]$ by (a), and $T=O$.

For (d), the discriminant is $\Delta=64n^6$. At an odd prime $p$, $v_p(\Delta)=6v_p(n)\le6<12$, so the
model is minimal at $p$. If $n$ is odd, then $v_2(\Delta)=6$ and the same holds at $2$. If $n=2m$ with
$m$ odd, then replacing $m$ by $mu^2$ with $u\in\Z_2^\times$ gives an isomorphic model over $\Z_2$, via
$(x,y)\mapsto(u^2x,u^3y)$. So minimality at $2$ depends only on $m$ modulo $(\Z_2^\times)^2$, that is, on
$m\bmod8$. Tate's algorithm (we used \texttt{ellminimalmodel} in PARI/GP~\cite{PARI}) shows that the model
is minimal at $2$ for $n=2,6,10,14$, which covers $m\equiv1,3,5,7\pmod8$.
\end{proof}

Now let $C$ be an elliptic curve over $\Q$ given by a Weierstrass equation
\begin{equation}\label{eq:weier}
y^2+a_1xy+a_3y=x^3+a_2x^2+a_4x+a_6,\qquad a_1,a_2,a_3,a_4,a_6\in\Z,
\end{equation}
and for a prime $\ell$ let $v_\ell$ be the $\ell$-adic valuation on $\Q$. If $(x,y)\in C(\Q)$ and
$t=-v_\ell(x)>0$, then $t$ is even. Indeed, the right-hand side of~\eqref{eq:weier} has valuation $-3t$,
because $x^3$ does and the other terms have valuation at least $-2t$. If $v_\ell(y)\ge-t$, every term on
the left has valuation at least $-2t$, which is impossible. So $v_\ell(y)<-t$. Then $y^2$ has strictly
smaller valuation than $a_1xy$ and $a_3y$, and comparing the two sides gives $2v_\ell(y)=-3t$.
Consequently, for $P\in C(\Q)\setminus\{O\}$ we can write $x(P)=A(P)/B(P)^2$ with $A(P),B(P)\in\Z$
coprime and $B(P)\ge1$, and for every prime $\ell$
\begin{equation}\label{eq:val}
\ell\mid B(P)\iff v_\ell\bigl(x(P)\bigr)<0\iff v_\ell\bigl(x(P)\bigr)=-2v_\ell\bigl(B(P)\bigr)\le-2.
\end{equation}
Since $-P$ has the same $x$-coordinate as $P$, $B(-P)=B(P)$. For a non-torsion point $Q$, the integers
$B_m=B(mQ)$, $m\ge1$, form the \emph{elliptic divisibility sequence} of $Q$~\cite{Ward}. A prime $\ell$ is
a \emph{primitive divisor} of $B_M$ if $\ell\mid B_M$ and $\ell\nmid B_m$ for $1\le m<M$.
Silverman~\cite{Si88} showed that all but finitely many terms have a primitive divisor. For the curves
$W_D\colon y^2=x^3+Dx$, where $D\ne0$ is an integer, there is a uniform statement. Note that
$E_n=W_{-n^2}$.

\begin{theorem}[Voutier--Yabuta~\cite{VY}, Verzobio~\cite{Ve}]\label{thm:VY}
Let $D$ be a fourth-power-free integer for which $y^2=x^3+Dx$ is a minimal model, and let $Q$ be a
non-torsion rational point on it with $x(Q)$ a square in $\Q$. Then $B_M(Q)$ has a primitive divisor for
every $M\ge3$.
\end{theorem}

Voutier and Yabuta~\cite{VY} prove that $B_M$ has a primitive divisor for every even $M>2$, and for every
odd $M>1$ when $x(Q)<0$ or $x(Q)$ is a square. Their proof needs the case $M=5m$ (their Lemma~5.1).
Verzobio~\cite{Ve} found a gap in the proof of that lemma and proved, for minimal models $y^2=x^3+Dx$, that
$B_{5m}$ has a primitive divisor for every $m$.

\section{Multiples of one point are never in arithmetic progression}

\begin{lemma}\label{lem:noAP}
Let $C$ be an elliptic curve over $\Q$ given by a Weierstrass equation~\eqref{eq:weier} with integer
coefficients, and let $Q\in C(\Q)$ be a non-torsion point such that $B_M(Q)$ has a primitive divisor for
every $M\ge3$. Let $m_{-1},m_0,m_1$ be nonzero integers whose absolute values are pairwise distinct. Then
\[
x(m_{-1}Q)+x(m_1Q)\ne2\,x(m_0Q).
\]
\end{lemma}

By Theorem~\ref{thm:VY}, the hypothesis on $Q$ holds when $C$ is $y^2=x^3+Dx$ with $D$ fourth-power-free,
the model is minimal, and $x(Q)$ is a square in $\Q$.

\begin{proof}
The three integers $|m_\alpha|$ are distinct and positive, so their maximum $M$ satisfies $M\ge3$. Let
$\ell$ be a primitive divisor of $B_M$, and let $\alpha^*$ be the index with $|m_{\alpha^*}|=M$. Since
$x(mQ)=x(|m|Q)$, \eqref{eq:val} gives
\[
v_\ell\bigl(x(m_{\alpha^*}Q)\bigr)=-2v_\ell(B_M)\le-2,
\qquad
v_\ell\bigl(x(m_\beta Q)\bigr)\ge0\quad(\beta\ne\alpha^*),
\]
because $\ell\nmid B_{|m_\beta|}$ when $|m_\beta|<M$.

If $\alpha^*=\pm1$, the left-hand side $x(m_{-1}Q)+x(m_1Q)$ has $\ell$-adic valuation at most $-2$. The
right-hand side $2x(m_0Q)$ has valuation at least $0$.

If $\alpha^*=0$, the right-hand side has valuation $v_\ell(2)-2v_\ell(B_M)\le1-2<0$. The left-hand side
has valuation at least $0$.

In either case the two sides are different.
\end{proof}

\section{Proof of Theorem~\ref{thm:main} and Corollary~\ref{cor:rank}}

\begin{proof}[Proof of Theorem~\ref{thm:main}]
Let $n=\sqf(b)$ and let $P_{-1},P_0,P_1\in E_n(\Q)$ be the points of~\eqref{eq:points}.
\begin{itemize}
\item By Lemma~\ref{lem:facts}(b), each $P_\alpha$ lies in $2E_n(\Q)$, since $x_\alpha$ and
$x_\alpha\pm n$ are squares.
\item The $x_\alpha=(E+\alpha a)/k^2$ are pairwise distinct because $a\neq0$, and none of the $P_\alpha$
is $O$.
\item Let $\Gamma=\langle P_{-1},P_0,P_1\rangle\subseteq2E_n(\Q)$. By Lemma~\ref{lem:facts}(c), $\Gamma$ is
a nonzero, torsion-free, finitely generated abelian group, so $\rk\Gamma\ge1$.
\end{itemize}
Suppose $\rk\Gamma=1$. Then $\Gamma$ is infinite cyclic, $\Gamma=\langle Q\rangle$, with
$Q\in\Gamma\subseteq2E_n(\Q)$ non-torsion, so $x(Q)$ is a square by Lemma~\ref{lem:facts}(b). Write
$P_\alpha=m_\alpha Q$ with $m_\alpha\in\Z\setminus\{0\}$. For a non-torsion $Q$, $x(mQ)=x(m'Q)$ implies
$mQ=\pm m'Q$, and hence $m=\pm m'$. So the distinctness of the $x_\alpha$ forces the $|m_\alpha|$ to be
pairwise distinct. Since $D=-n^2$ is fourth-power-free and the model is minimal
(Lemma~\ref{lem:facts}(d)), Theorem~\ref{thm:VY} shows that $B_M(Q)$ has a primitive divisor for every
$M\ge3$. So Lemma~\ref{lem:noAP} applies, and it contradicts $x_{-1}+x_1=2x_0$. Hence $\rk\Gamma\ge2$.

The statement for the $a$-progressions follows by exchanging the roles of $a$ and $b$, which needs
$b\neq0$.
\end{proof}

\begin{proof}[Proof of Corollary~\ref{cor:rank}]
The common differences of the two diagonals of~\eqref{eq:lucas} are $a$ and $b$, and $ab\ne0$ for a
magic square of squares. By Theorem~\ref{thm:main}, $\rk E_n(\Q)\ge2$ for $n=\sqf(a)$ and $n=\sqf(b)$; in
particular $n$ is a congruent number.

For the density, $E_n$ is the quadratic twist by $n$ of $E_1\colon y^2=x^3-x$. The curve $E_1$ has full
rational $2$-torsion. It has no rational cyclic subgroup of order four, because the $x$-coordinates of
its points of exact order four are the roots of $x^2+1$ and $x^2\pm2x-1$, none of which is rational.
Smith~\cite{Smith} proved that for such a curve the number of twists $E_1^d$ with $|d|<N$ and rank at
least two is $o(N)$.

For the list, PARI/GP's \texttt{ellrank}~\cite{PARI} proves $\rk E_n(\Q)\ge2$ for $n=34,41,65,137,138,145$,
and $\rk E_n(\Q)\le1$ for every other squarefree $n\le150$ except $n=113$, where it gives
$0\le\rk\le2$. For $n=113$, Tunnell's criterion~\cite{Tunnell} settles the question. If an odd
squarefree $n$ is congruent then
$\#\{x^2+2y^2+8z^2=n\}=2\,\#\{x^2+2y^2+32z^2=n\}$. For $n=113$ these counts are $32$ and $24$, so $113$
is not congruent, and $\rk E_{113}(\Q)=0$.
\end{proof}

\section{Remarks}

\begin{remark}[Relation to~\cite{BST}]\label{rem:BST}
Every magic square of squares is a square multiple of a primitive one, whose entries have no common factor
greater than $1$, and the property in Theorem~\ref{thm:main} is unchanged by scaling. In a primitive one,
write the center as $E=e^2$. Each line through the center is $(e^2-t,e^2,e^2+t)$ with
$t\in\{a,b,a+b,a-b\}$, and $e^2-t=x^2$, $e^2+t=y^2$ give $x^2+y^2=2e^2$.
\begin{itemize}
\item If $e$ were even, then $x^2+y^2\equiv0\pmod4$ would force $x$ and $y$ to be even for every such
line, making all nine entries even. So $e$ is odd.
\item Then $x^2+y^2=2e^2\equiv2\pmod4$ forces $x$ and $y$ odd, so all nine entries are odd and
$t=y^2-e^2\equiv0\pmod8$. In particular $8\mid a$ and $8\mid b$.
\end{itemize}
Writing $|b|=nk^2$ with $n$ squarefree, $v_2(n)\le1$ forces $k$ to be even. So
$x_\alpha=(E+\alpha a)/k^2$ has an odd numerator and an even denominator, and is never an integer.
Theorem~1.1 of~\cite{BST}, which concerns integral progressions on $E_n$ with $n$ squarefree, therefore
never applies to the points~\eqref{eq:points}; nor does their Theorem~1.4 on large squarefree twists.
Conversely, Lemma~\ref{lem:noAP} and Theorem~\ref{thm:VY} do not recover~\cite[Thm.~1.1]{BST}. That
theorem allows torsion points, and points outside $2E_n(\Q)$, to which Theorem~\ref{thm:VY} need not
apply.
\end{remark}

\begin{remark}[What the argument uses]\label{rem:uses}
The proof uses that all three $b$-progressions (or all three $a$-progressions) consist of squares. Since
these partition the nine entries, it uses all nine square conditions. It says nothing about magic squares
with seven or eight square entries. For instance, Bremner's magic square with seven square
entries~\cite{Bremner99} has $E=425^2$, $a=-41496$ and $b=138600$. Two of its $b$-progressions and one of
its $a$-progressions consist of squares, and $\sqf(b)=154$ while $\sqf(a)=10374$. Here $\rk E_{154}(\Q)=2$
and $\rk E_{10374}(\Q)=1$, which is consistent with the theorem, since that square is not a magic square
of squares.
\end{remark}

\begin{remark}[Near misses]\label{rem:near}
Two full $b$-progressions already give two of the points~\eqref{eq:points}, without the third
$b$-progression and without a square center. These two points need not be independent, even in a magic
square with nine distinct entries. For $E=1201^2$, $a=1319901$ and $b=117600=6\cdot140^2$, the
square~\eqref{eq:lucas} is
\[
\begin{pmatrix}
2762302 & 70^2 & 1249^2\\
490^2 & 1201^2 & 2644702\\
1151^2 & 2879902 & 350^2
\end{pmatrix},
\]
with magic sum $4327203$ and six square entries. Its $b$-progressions with $\alpha=-1$ and $\alpha=0$ are
full, and the first of them is $70^2\cdot(1,25,49)$. With $G=(-3,9)$ and $H=2G$ as in
Example~\ref{ex:6}, we have $x_{-1}=25/4=x(H)$ and $x_0=(1201/140)^2=x(2H)$, so $P_{-1}=\pm H$ and
$P_0=\pm2H$. Since $\rk E_6(\Q)=1$, the hypothesis of Theorem~\ref{thm:main} cannot be weakened to two
full $b$-progressions.

In the near misses we have checked, however, the two points are independent. These are Bremner's square
of Remark~\ref{rem:uses}, whose $b$-progressions with $\alpha=0,1$ are full, and $2{,}320$ configurations,
found by a computer search over centers $E<10^{12}$, in which the $b$-progressions with $\alpha=\pm1$ are
full and $E$ is not a square. The proof is exact. Write the two points as $P=2R$ and $P'=2R'$, as
Lemma~\ref{lem:facts}(b) allows. The map $(x,y)\mapsto(x,x-n)$, with $(0,0)\mapsto(-1,-n)$ and
$(n,0)\mapsto(n,2)$, induces an injective homomorphism
$\phi\colon E_n(\Q)/2E_n(\Q)\to(\Q^\times/\Q^{\times2})^2$~\cite[Prop.~X.1.4]{AEC}. In all but one case,
$\phi$ maps $R$, $R'$, $(0,0)$ and $(n,0)$ to elements that are linearly independent over $\mathbb{F}_2$.
Then $R$ and $R'$ are independent modulo $E_n(\Q)_{\mathrm{tors}}$, which is $E_n[2]$ by
Lemma~\ref{lem:facts}(a): a relation $cR+c'R'\in E_n[2]$ with $c,c'$ coprime would give a nontrivial
relation modulo $2E_n(\Q)$. Hence $P$ and $P'$ are independent. The exception is a configuration with
$n=210$, in which $\phi$ maps $R$, $R'$ and $(n,0)$ to elements with product~$1$. Then
$R+R'+(n,0)=2S$ for some $S\in E_n(\Q)$, the same test succeeds for $S$ and $R'$, and $R\equiv2S-R'$
modulo torsion. So in all these near misses $\rk E_{\sqf(b)}(\Q)\ge2$,
and the rank bound of Corollary~\ref{cor:rank} does not by itself separate magic squares of squares from
them. The accompanying script \texttt{near\_misses.gp} carries out these computations.
\end{remark}

\begin{remark}[Rank at least two]
The valuation argument of Lemma~\ref{lem:noAP} shows, in any rank, that a prime dividing the denominator
of one of $x_{-1}$, $x_0$, $x_1$ must divide the denominator of another. When $\Gamma$ has rank at least two,
this becomes a condition on the common prime factors of the denominators of points in a lattice of rank
at least two, that is, of the terms of an elliptic net~\cite{Stange}. Bounds of the required kind for the
greatest common divisor of $B(P)$ and $B(Q)$, for independent $P$ and $Q$, follow from Vojta's
conjecture~\cite{Si05}. They are not known unconditionally, let alone uniformly in $n$. So we do not
know how to extend Theorem~\ref{thm:main} to rank two.
\end{remark}

\begin{remark}[Computations]
As a check on Lemma~\ref{lem:noAP}, consider the $116$ squarefree $n\le400$ for which PARI/GP, run with
the random seed of the accompanying script \texttt{checks.gp}, proves $\rk E_n(\Q)=1$ and finds a
generator $G$, and set $H=2G$. For ten more $n$, such as $n=157$, it proves rank one but finds no
generator. In each case, no three of $x(H),\dots,x(16H)$ are in arithmetic progression, and every $B_m(H)$
with $2\le m\le16$ has a primitive divisor. The same script reproduces the other computations in this
note, except those of Remark~\ref{rem:near}.
\end{remark}

\begin{remark}[Formal verification]\label{rem:lean}
Lemma~\ref{lem:noAP}, Theorem~\ref{thm:main} and the rank statement of Corollary~\ref{cor:rank} are
formalized in Lean~4~\cite{Lean} with the Mathlib library~\cite{Mathlib}, in the file
\texttt{RankOne.lean} that accompanies this note. Two inputs enter as explicit hypotheses: the conclusion of
Theorem~\ref{thm:VY} for $E_n$, and Lemma~\ref{lem:facts}(a). Everything else is proved, including the
halving formula behind Lemma~\ref{lem:facts}(b), the group-law computations and the passage from a magic
square of squares to the points~\eqref{eq:points}.
\begin{itemize}
\item Lemma~\ref{lem:noAP} is proved as stated, for every Weierstrass equation~\eqref{eq:weier} with
integer coefficients, together with the facts~\eqref{eq:val} about valuations that its proof uses.
\item The rank of $\Gamma$ is its rank as a free $\Z$-module.
\item $\rk E_n(\Q)\ge2$ is stated as the existence of two $\Z$-linearly independent points of $E_n(\Q)$.
\end{itemize}
Lean reports that the proofs use only its three standard axioms. The density statement and the list in
Corollary~\ref{cor:rank} are not formalized.
\end{remark}

\subsection*{Acknowledgments}
The argument, the literature checks, the computations and the Lean formalization were developed with the
assistance of Claude (Anthropic). Computations used PARI/GP.

\begin{thebibliography}{99}

\bibitem{AS} A.~Auel and B.~Singer, \emph{The algebraic geometry of 3-by-3 magic squares of squares},
arXiv:2609.09351, 2026.

\bibitem{Boyer} C.~Boyer, \emph{Some notes on the magic squares of squares problem}, Math. Intelligencer
\textbf{27} (2005), no.~2, 52--64.

\bibitem{Bremner99} A.~Bremner, \emph{On squares of squares}, Acta Arith. \textbf{88} (1999), 289--297.

\bibitem{Bremner01} A.~Bremner, \emph{On squares of squares II}, Acta Arith. \textbf{99} (2001), 289--308.

\bibitem{BST} A.~Bremner, J.~H.~Silverman and N.~Tzanakis, \emph{Integral points in arithmetic
progression on $y^2=x(x^2-n^2)$}, J. Number Theory \textbf{80} (2000), 187--208.

\bibitem{BTVA} N.~Bruin, J.~Thomas and A.~V\'arilly-Alvarado, \emph{Explicit computation of symmetric
differentials and its application to quasihyperbolicity}, Algebra Number Theory \textbf{16} (2022),
no.~6, 1377--1405.

\bibitem{Choi} S.~Choi, \emph{Additive rigidity for $x$-coordinates of rational points on elliptic
curves}, arXiv:2510.03828, 2025 (revised 2026).

\bibitem{Kn} A.~W.~Knapp, \emph{Elliptic Curves}, Mathematical Notes \textbf{40}, Princeton University
Press, 1992.

\bibitem{Ko} N.~Koblitz, \emph{Introduction to Elliptic Curves and Modular Forms}, 2nd ed., Graduate Texts
in Mathematics \textbf{97}, Springer, 1993.

\bibitem{Kominers} S.~D.~Kominers, \emph{Center-zero magic squares of squares over number fields},
preprint, 2026,
\url{https://www.scottkom.com/assets/articles/Kominers_Center-Zero_Magic_Squares.pdf}.

\bibitem{LaBar} M.~LaBar, \emph{Problem 270}, College Math. J. \textbf{15} (1984), 69.

\bibitem{Mathlib} The mathlib Community, \emph{The Lean mathematical library}, in: Proceedings of the 9th
ACM SIGPLAN International Conference on Certified Programs and Proofs (CPP 2020), ACM, 2020, 367--381.

\bibitem{Lean} L.~de~Moura and S.~Ullrich, \emph{The Lean~4 theorem prover and programming language}, in:
Automated Deduction -- CADE~28, Lecture Notes in Comput. Sci. \textbf{12699}, Springer, 2021, 625--635.

\bibitem{PARI} The PARI~Group, \emph{PARI/GP version 2.17.3}, Univ. Bordeaux,
\url{https://pari.math.u-bordeaux.fr/}.

\bibitem{Robertson} J.~P.~Robertson, \emph{Magic squares of squares}, Math. Mag. \textbf{69} (1996),
no.~4, 289--293.

\bibitem{Si88} J.~H.~Silverman, \emph{Wieferich's criterion and the $abc$-conjecture}, J. Number Theory
\textbf{30} (1988), 226--237.

\bibitem{Si05} J.~H.~Silverman, \emph{Generalized greatest common divisors, divisibility sequences, and
Vojta's conjecture for blowups}, Monatsh. Math. \textbf{145} (2005), 333--350; arXiv:math/0407415.

\bibitem{AEC} J.~H.~Silverman, \emph{The Arithmetic of Elliptic Curves}, 2nd ed., Graduate Texts in
Mathematics \textbf{106}, Springer, 2009.

\bibitem{Smith} A.~Smith, \emph{$2^\infty$-Selmer groups, $2^\infty$-class groups, and Goldfeld's
conjecture}, arXiv:1702.02325, 2017.

\bibitem{Stange} K.~E.~Stange, \emph{Elliptic nets and elliptic curves}, Algebra Number Theory \textbf{5}
(2011), no.~2, 197--229.

\bibitem{Tunnell} J.~B.~Tunnell, \emph{A classical Diophantine problem and modular forms of weight $3/2$},
Invent. Math. \textbf{72} (1983), 323--334.

\bibitem{Ve} M.~Verzobio, \emph{Primitive divisors of elliptic divisibility sequences for elliptic curves
with $j=1728$}, Acta Arith. \textbf{198} (2021); arXiv:2001.09634.

\bibitem{VY} P.~Voutier and M.~Yabuta, \emph{Primitive divisors of certain elliptic divisibility
sequences}, Acta Arith. \textbf{151} (2012), 165--190.

\bibitem{Ward} M.~Ward, \emph{Memoir on elliptic divisibility sequences}, Amer. J. Math. \textbf{70}
(1948), 31--74.

\end{thebibliography}

\end{document}
